%% 6-0 序章　「次の瞬間を決めるルール」で未来を予想する（第14回：組 V6-1。ページ制限なし。偶数ページ）
%% ① 次の状態＝今の状態＋変わる速さ×短い時間（2-18 の冷めるお湯・5-8 の1行）② 天気予報（はてな？の答え）と、その歴史
%% ③ ロケット・人工衛星・ゲームのボール ④ ルールとはじめの状態 ⑤ 3つの部と「2周読む」 ⑥ ばね・熱・電気回路
%% ⑦ 乗りかえ駅 ⑧ 新しい記号と、この巻の読み方
\makeatletter
\gdef\bk@partlabel{第6巻}\gdef\bk@partname{未来を計算する}
%% ---- この章だけで使う部品（頭文字 \bkvia…） ----
% ばね（ジグザグ）  \bkviaspring{(左はし)}{長さ mm}{山の数}{山の高さ mm}{色}
\newcommand\bkviaspring[5]{\begin{scope}[shift={#1}]
  \pgfmathsetmacro\bkvia@w{#2/(#3*2+2)}
  \draw[#5, line width=0.9pt, line join=round] (0,0) -- (\bkvia@w,0)
    \foreach \q in {1,...,#3} { -- ({\bkvia@w*(2*\q-0.5)},#4) -- ({\bkvia@w*(2*\q+0.5)},-#4)} -- ({#2-\bkvia@w},0) -- (#2,0);
  \end{scope}}
% かべ（ななめの線つき）  \bkviawall{(下のはし)}{高さ}
\newcommand\bkviawall[2]{\begin{scope}[shift={#1}]
  \fill[bkGray!25] (-3,0) rectangle (0,#2);
  \draw[bkGray!80!black, line width=0.9pt] (0,0) -- (0,#2);
  \end{scope}}
% 路線の1行  \bkviaroute{y}{路線番号}{前の巻}{この巻の駅}{この先}
\newcommand\bkviaroute[5]{%
  \draw[bkRoute#2, line width=3pt, line cap=round] (10,#1) -- (132,#1);
  \node[inner sep=0pt] at (4,#1) {\bkrouteicon[7mm]{#2}};
  \fill[white] (24,#1) circle (2.2);\draw[bkRoute#2, line width=1.1pt] (24,#1) circle (2.2);
  \node[anchor=north, font=\sffamily\footnotesize, text=bkGray, inner sep=0pt] at (24,#1-3.8) {#3};
  \fill[white, rounded corners=1.5mm] (38,#1-4.4) rectangle (110,#1+4.4);
  \draw[bkRoute#2, line width=1.6pt, rounded corners=1.5mm] (38,#1-4.4) rectangle (110,#1+4.4);
  \node[font=\sffamily\footnotesize, text=bkInk] at (74,#1) {#4};
  \fill[white] (124,#1) circle (2.2);\draw[bkRoute#2, line width=1.1pt] (124,#1) circle (2.2);
  \node[anchor=north, font=\sffamily\footnotesize, text=bkGray, inner sep=0pt] at (124,#1-3.8) {#5};}
% 部の地図の1まい  \bkviapart{上のはしの y}{部の名前}{章}{絵}{中身}   はば132mm・高さ40mm
\newcommand\bkviapart[5]{%
  \draw[bkTheme!80!black, line width=1pt, fill=white, rounded corners=2.5mm] (0,#1) rectangle (132,#1-38);
  \node[anchor=north west, fill=bkTheme, text=bkOn, rounded corners=1.2mm, font=\sffamily\bfseries\small,
    inner xsep=2mm, inner ysep=1mm] at (3.5,#1-3.5) {#2};
  \node[anchor=north west, font=\sffamily\footnotesize, text=bkGray, inner sep=0pt] at (4.5,#1-11.5) {#3};
  \begin{scope}[shift={(23,#1-26)}]#4\end{scope}
  \node[anchor=north west, text width=78mm, align=left, font=\small, inner sep=0pt] at (51,#1-4.5) {#5};}
% 箇条書きの1項目  \bkviaitem{文}
\newcommand\bkviaitem[1]{\par\noindent\hangindent=1\zw\hangafter=1 ・#1\par}
% 会話の1行  \bkviatalk{y}{sora|mina|delta}{名前}{せりふ}
\newcommand\bkviatalk[4]{%
  \bkChar{#2}{(7,#1-6.5)}{12mm}
  \node[rectangle callout, callout absolute pointer={(15.5,#1)}, anchor=west, draw=bkTheme!80!black, line width=0.8pt,
    fill=white, rounded corners=2.5mm, inner xsep=3mm, inner ysep=1.8mm, text width=103mm, align=left,
    font=\small] at (21,#1) {#4};
  \node[anchor=north, font=\sffamily\footnotesize, text=bkGray, inner sep=0pt] at (7,#1-7.5) {#3};}
% 読み方のカード  \bkviacard{(x,y)}{アイコン}{色}{見出し}{中身}  はば64mm・高さ50mm
\newcommand\bkviacard[5]{%
  \node[anchor=north west, draw=#3, line width=1pt, fill=white, rounded corners=2.5mm,
    minimum width=64mm, minimum height=54mm, inner sep=0pt] at #1 {};
  \node[anchor=north west, inner sep=0pt] (ci) at ($#1+(3.5,-3)$) {\bkicon[1.5em]{#2}{#3}};
  \node[anchor=west, font=\sffamily\bfseries\small, text=#3!85!black, inner sep=0pt] at ([xshift=2mm]ci.east) {#4};
  \node[anchor=north west, text width=57mm, align=left, font=\small, inner sep=0pt] at ($#1+(3.5,-12)$) {#5};}
% 見出しのアイコン（路線：線と2つの駅。明るい青緑の地なので線は濃い紺）
\newcommand\bk@icon@bkviaroute{%
  \draw[bkOn, line width=0.13*\bk@iu, line cap=round] (-0.62,0) -- (0.62,0);
  \filldraw[fill=white, draw=bkOn, line width=0.09*\bk@iu] (-0.36,0) circle (0.2) (0.36,0) circle (0.2);}
% 太陽・雲・雨の小さな絵  \bkviasun{(x,y)}{半径}  \bkviacloud{(x,y)}{大きさ}  \bkviarain{(x,y)}{大きさ}
\newcommand\bkviasun[2]{\begin{scope}[shift={#1}]
  \foreach \a in {0,45,...,315} {\draw[bkOrange, line width=0.9pt, line cap=round] (\a:{#2*1.35}) -- (\a:{#2*1.8});}
  \filldraw[fill=bkYellow, draw=bkOrange, line width=0.8pt] (0,0) circle (#2);\end{scope}}
\newcommand\bkviacloud[2]{\begin{scope}[shift={#1}, scale=#2]
  \filldraw[fill=white, draw=bkGray!80!black, line width=0.7pt] (-3.2,-1.6) arc (270:90:1.6) arc (180:30:2.2) arc (120:0:1.7) arc (90:-90:1.6) -- cycle;
  \end{scope}}
\newcommand\bkviarain[2]{\begin{scope}[shift={#1}, scale=#2]
  \filldraw[fill=bkGray!25, draw=bkGray!80!black, line width=0.7pt] (-3.2,-1.6) arc (270:90:1.6) arc (180:30:2.2) arc (120:0:1.7) arc (90:-90:1.6) -- cycle;
  \foreach \x in {-2.2,0,2.2} {\draw[bkBlue, line width=0.8pt, line cap=round] (\x,-2.6) -- (\x-0.6,-4.2);}
  \end{scope}}
% 空気の箱の位置（i 右へ、j 奥へ、k 上へ）
\def\bkviaP(#1,#2,#3){({10*(#1)+5*(#2)},{3.6*(#2)+8*(#3)})}
\makeatother

\bkchapter{6}{0}{「次の瞬間を決めるルール」で\\未来を予想する}{\bklevel{2}}{「次の瞬間を決めるルール」で未来を予想する}

\noindent テレビやスマホの天気予報は、「明日の午後は雨。最高気温は18\,℃」のように、まだ来ていない明日のことを教えてくれます。週間予報では、1週間先の天気まで出ています。

\begin{bkfigure}
\begin{tikzpicture}[x=1mm, y=1mm]
  % スマホの画面（3日分の予報）
  \filldraw[fill=bkInk!85, draw=bkInk, line width=0.8pt, rounded corners=4mm] (0,0) rectangle (112,46);
  \fill[white, rounded corners=2mm] (4,4) rectangle (108,42);
  \node[anchor=north west, font=\sffamily\bfseries\footnotesize, text=bkInk] at (8,40) {天気予報};
  \draw[bkLine, line width=0.5pt] (8,33.5) -- (104,33.5);
  \foreach \i/\d/\t/\p in {0/今日/22/10, 1/明日/18/80, 2/あさって/20/30} {
    \node[font=\sffamily\bfseries\footnotesize, text=bkInk] at ({22+34*\i},30) {\d};
    \node[font=\sffamily\footnotesize, text=bkInk] at ({22+34*\i},10.5) {最高 \t\,℃};
    \node[font=\sffamily\footnotesize, text=bkGray] at ({22+34*\i},6.2) {雨の確率 \p\,\%};}
  \bkviasun{(22,19.5)}{3.6}
  \bkviarain{(56,20.5)}{1.5}
  \bkviasun{(87,22.5)}{2.6}
  \bkviacloud{(93,18)}{1.3}
  \draw[bkLine, line width=0.5pt] (39,8) -- (39,32) (73,8) -- (73,32);
  % 説明
  \node[anchor=west, text width=16mm, align=left, font=\sffamily\footnotesize, text=bkDeep] at (115,23) {明日も\\あさっても\\わかる？};
\end{tikzpicture}
\bkfigcaption{スマホの天気予報。明日やあさっての天気まで出ている。}
\end{bkfigure}

明日の天気は、まだだれも見ていません。それなのに、天気予報はかなりよく当たります。気象庁の発表では、明日の「雨が降るか降らないか」の予報は、8割以上当たっています。どうして、まだ起きていないことがわかるのでしょう。

\begin{hatena}
\par\addvspace{4pt}\noindent\begin{minipage}[c]{0.56\linewidth}
天気予報は、まだ来ていない明日のことを、なぜ計算できるのかな？\par\vspace{3pt}
ロケットや人工衛星の通り道も、何年も先まで計算されています。同じ考え方なのかな？
\end{minipage}\hfill
\begin{minipage}[c]{0.4\linewidth}\centering
\begin{tikzpicture}[x=1mm, y=1mm]
  \node[draw=bkGray!80!black, line width=0.8pt, fill=white, rounded corners=1.5mm, minimum width=15mm, minimum height=10mm, font=\sffamily\footnotesize] (a) at (0,0) {今日};
  \node[draw=bkGray!80!black, line width=0.8pt, fill=white, rounded corners=1.5mm, minimum width=15mm, minimum height=10mm, font=\sffamily\footnotesize] (b) at (36,0) {明日};
  \draw[bk link] (a) -- node[above, font=\sffamily\bfseries\large, text=bkOrange!80!black, inner sep=1.5mm] {？} (b);
  \bkviasun{(0,11)}{2.2}
  \bkviarain{(36,12.5)}{0.9}
\end{tikzpicture}
\end{minipage}
\end{hatena}

答えのかぎは、この巻の題名にもなっている「\textbf{次の瞬間を決めるルール}」です。序章では、このルールで未来を予想するしくみを、冷めるお湯・天気予報・ボールの3つの場面で見ていきます。そのあとで、この巻の地図を作ります。

\begin{omoidasou}
\bkrecall{2-18}{冷めるお湯の表：部屋は20\,℃。1分ごとに、部屋との差の5\,\%ずつ温度が下がる。}
\bkrecall{5-8}{次の位置 $\approx$ 今の位置＋今の速度×$\Delta t$。$\Delta t$ が短いほど、予想はよく当たる（オイラー法）。}
\bkrecall{5-13}{法則とはじめの状態がそろって、はじめて先のようすが1つに決まる。}
\bkrecall{5-18}{答えの候補リスト。冷めるお湯の温度は $T=20+60e^{-0.05t}$（はじめ80\,℃）。}
\end{omoidasou}

\begin{yattemiyou}[80\,℃のお湯の、少し先を予想しよう]
部屋の温度は20\,℃です。カップのお湯は、「\textbf{1分間に、部屋との温度の差の5\,\%だけ下がる}」というルールで冷めていきます（2-18と同じルール）。いま、お湯は80\,℃です。電卓を使って、未来の温度を予想しましょう。
\begin{enumerate}[label={\bkcircled[bkGreen]{\arabic*}}, leftmargin=2.4\zw, itemsep=1pt, topsep=3pt]
\item いまの、部屋との差は何\,℃かな？\enspace その5\,\%は何\,℃かな？
\item 1分後の温度は何\,℃かな？
\item 1分後の温度を「いま」として、同じ計算をくり返そう。2分後、3分後……5分後まで、表にしよう。
\end{enumerate}
\begin{bkcenter}
\begin{bktable}{colspec={Q[l,m,wd=24mm] *{6}{Q[c,m,wd=11mm]}}, rowsep=1.5pt}
時刻（分） & 0 & 1 & 2 & 3 & 4 & 5 \\
温度（℃） & 80 & ？ & ？ & ？ & ？ & ？ \\
\end{bktable}
\end{bkcenter}
\end{yattemiyou}

\noindent ①は、$80-20=60$ で、差は60\,℃です。その5\,\%は $60\times0.05=3$ で、3\,℃です（5\,\%は $\frac{5}{100}=0.05$ 倍）。②は、$80-3=77$ で、1分後は77\,℃です。③は、77\,℃を「いま」にして、同じ計算をします。差は $77-20=57$、その5\,\%は $57\times0.05=2.85$、だから2分後は $77-2.85=74.15$ です。これをくり返すと、次の表になります（小数第3位で四捨五入）。

\begin{bkcenter}
\begin{bktable}{colspec={Q[l,m,wd=37mm] *{6}{Q[c,m,wd=11.5mm]}}, rowsep=1.5pt, colsep=3pt}
時刻（分） & 0 & 1 & 2 & 3 & 4 & 5 \\
温度（℃） & 80 & 77 & 74.15 & 71.44 & 68.87 & 66.43 \\
部屋との差（℃） & 60 & 57 & 54.15 & 51.44 & 48.87 & 46.43 \\
1分で下がる温度（℃） & 3 & 2.85 & 2.71 & 2.57 & 2.44 & 2.32 \\
\end{bktable}
\bktabcaption{1分ずつ進めた、お湯の温度の予想}
\end{bkcenter}

同じ計算を10分後まで続けると、約55.92\,℃になります。まだ5分も10分もたっていないのに、計算だけで「5分後」「10分後」の温度がわかりました。これが、未来を計算するということです。

%% ---------------- ① 次の瞬間を決めるルール ----------------
\bksection{doguu}{次の瞬間を決めるルール}

\bksubsection{1}{今の状態から、少し先の状態へ}

\label{bk:6-0:rule}やってみようで、毎回したことは、次の2つだけでした。

\begin{enumerate}[label={\bkcircled[\bkTT]{\arabic*}}, leftmargin=2.4\zw, itemsep=2pt, topsep=3pt]
\item \textbf{いまの温度から、下がる速さを求める。}\enspace 「1分あたり、部屋との差の5\,\%」というルールを使う。80\,℃なら「1分あたり3\,℃」。
\item \textbf{下がる速さに時間をかけて、いまの温度から引く。}\enspace 1分なら $3\times1=3$ で、$80-3=77$。
\end{enumerate}

①で使った「いまの状態を見て、どれだけの速さで変わるかを教えてくれるきまり」を、この本では\textbf{次の瞬間を決めるルール}とよびます。\bkidx{つぎのしゅんかんをきめるるーる}{次の瞬間を決めるルール}②は、5-8の「次の位置 $\approx$ 今の位置＋今の速度×$\Delta t$」と同じ形の計算です。まとめると、次のようになります。

\begin{taisetsu}
\centering\textbf{次の状態}\enspace ＝\enspace\textbf{今の状態}\enspace ＋\enspace\textbf{変わる速さ}\enspace ×\enspace\textbf{短い時間}\par\vspace{3pt}
{\small 変わる速さは、ルールを使って「今の状態」から求める。}
\end{taisetsu}

温度が下がるときは、変わる速さをマイナスの数にすると、「たす」の形のままで書けます。80\,℃のときの変わる速さは「1分あたり $-3$\,℃」なので、$80+(-3)\times1=77$ です。この形は、下がるときも上がるときも同じ式で書けるので便利です。

この計算は、くり返すことができます。求めた「少し先の状態」を、新しい「今の状態」にして、もう一度ルールを使えばよいのです。図にすると、ぐるぐる回る輪になります。

\begin{bkfigure}
\begin{tikzpicture}[x=1mm, y=1mm,
  st/.style={bk node, minimum width=34mm, minimum height=15mm, font=\sffamily\small}]
  \node[st] (now) at (0,0) {\textbf{今の状態}\\{\footnotesize 温度 80\,℃}};
  \node[st] (sp) at (66,0) {\textbf{変わる速さ}\\{\footnotesize 1分あたり $-3$\,℃}};
  \node[st] (nx) at (66,-34) {\textbf{少し先の状態}\\{\footnotesize 1分後 77\,℃}};
  \draw[bk link] (now) -- node[bk link label, above] {ルールを使う} (sp);
  \draw[bk link] (sp) -- node[bk link label, right, align=left] {×1分を\\今の状態にたす} (nx);
  \draw[bk link, bkOrange!80!black] (nx.west) -| node[bk link label, pos=0.25, below, text=bkOrange!80!black] {これを「今」にして、くり返す} (now.south);
\end{tikzpicture}
\bkfigcaption{次の瞬間を決めるルールで、少しずつ未来へ進む}
\label{bk:6-0:fig-wa}
\end{bkfigure}

輪を1回まわるごとに、1分先へ進みます。10回まわれば10分後、60回まわれば1時間後の温度がわかります。電卓では大変でも、コンピュータなら、1秒に何億回でもまわせます。

\bksubsection{2}{ルールは「今」しか見ていない}

\label{bk:6-0:ima}ルールには、大切な性質があります。2つのカップでくらべてみましょう。

\begin{bkfigure}
\begin{tikzpicture}[x=1mm, y=1mm]
  % カップA・カップB
  \foreach \x/\lab/\txt in {0/A/{8分前に80\,℃のお湯を入れた。\\いま約60\,℃。}, 70/B/{いま、60\,℃のお湯を\\入れたばかり。}} {
    \begin{scope}[shift={(\x,0)}]
      \filldraw[fill=bkBlue!12, draw=bkInk!70, line width=0.9pt] (0,0) -- (-1.5,18) -- (17.5,18) -- (16,0) -- cycle;
      \draw[bkInk!70, line width=0.9pt] (17.1,13) arc (90:-90:3.6 and 4.2);
      \foreach \s in {3,8,13} {\draw[bkGray, line width=0.7pt] (\s,20) .. controls (\s-1.5,22.5) and (\s+1.5,24) .. (\s,26.5);}
      \node[font=\sffamily\bfseries\small] at (8,9) {\lab};
      \node[anchor=west, text width=36mm, align=left, font=\footnotesize] at (24,12) {\txt};
    \end{scope}}
  \node[anchor=north, text width=120mm, align=center, font=\sffamily\small, text=bkDeep] at (60,-4)
    {どちらも、次の1分で約2\,℃下がる（差40\,℃の5\,\%）};
\end{tikzpicture}
\bkfigcaption{いま同じ60\,℃なら、これまでのことにかかわらず、次の1分の下がり方は同じ}
\end{bkfigure}

カップAは、8分前に80\,℃だったお湯が冷めて、いま約60\,℃です。カップBは、いま60\,℃のお湯を入れたばかりです。これまでのことはちがいますが、次の1分で下がる温度は、どちらも $(60-20)\times0.05=2$ で、約2\,℃です。ルールが見ているのは「いまの温度」だけだからです。

この性質のおかげで、「今の状態」さえわかれば、そこから先を計算できます。昔のことを全部おぼえておく必要はありません。

\bksubsection{3}{時間を短く区切るほど、よく当たる}

\label{bk:6-0:mijikai}やってみようでは、1分ごとに区切りました。でも、ほんとうは1分のあいだにも、お湯は冷めて、下がる速さは小さくなっていきます。表の1行めで、0分では1分あたり3\,℃だった下がる速さが、1分後には2.85\,℃になっています。それなのに、1分間ずっと「3\,℃」の速さで下がったとして計算しました。だから、予想は少しずれます。

5-18で求めた式 $T=20+60e^{-0.05t}$ を使うと、10分後のほんとうの温度は約56.39\,℃です。1分ごとの予想55.92\,℃は、約0.47\,℃低くなっています。区切りを0.5分、0.1分と短くすると、予想はほんとうの値に近づきます。このことは、6-2でくわしく調べます。

\begin{bkfigure}
\begin{tikzpicture}
\begin{axis}[bkplot, width=104mm, height=62mm, xmin=0, xmax=11, ymin=50, ymax=84, xtick={0,2,...,10}, ytick={50,60,70,80},
  xlabel={時刻（分）}, ylabel={温度（℃）}]
  \addplot[bkcurve, domain=0:10.5, samples=60] {20+60*exp(-0.05*x)};
  \addplot[only marks, mark=*, mark size=1.7pt, bkOrange] coordinates {(0,80) (1,77) (2,74.15) (3,71.44) (4,68.87) (5,66.43) (6,64.11) (7,61.90) (8,59.81) (9,57.82) (10,55.92)};
  \node[anchor=south west, font=\sffamily\footnotesize, text=bkBlue] at (axis cs:4.6,70.5) {ほんとうの温度（5-18の式）};
  \node[anchor=north east, font=\sffamily\footnotesize, text=bkOrange!80!black, fill=white, inner sep=1pt] at (axis cs:4.2,65.5) {1分ごとの予想};
  \draw[bkGray, line width=0.4pt] (axis cs:4.3,65.2) -- (axis cs:4.9,67.6);
  \node[anchor=west, font=\sffamily\footnotesize, text=bkBlue] at (axis cs:10.3,57.8) {56.39};
  \node[anchor=west, font=\sffamily\footnotesize, text=bkOrange!80!black] at (axis cs:10.3,54.3) {55.92};
  \draw[bkGray, line width=0.4pt] (axis cs:10.1,56.3) -- (axis cs:10.35,57.6) (axis cs:10.1,55.85) -- (axis cs:10.35,54.6);
\end{axis}
\end{tikzpicture}
\bkfigcaption{1分ごとの予想（オレンジの点）と、ほんとうの温度（青の曲線）。少しずつ下にずれていく。}
\end{bkfigure}

%% ---------------- ② 天気予報 ----------------
\bksection{butsuri}{天気予報も、同じ考え方}

\bksubsection{4}{空気を箱に分けて、ルールで進める}

\label{bk:6-0:tenki}天気予報も、冷めるお湯と同じ考え方で計算しています。\bkidx{てんきよほう}{天気予報}ちがうのは、あつかう「状態」の数が、とても多いことです。

\begin{bkfigure}
\begin{tikzpicture}[x=1mm, y=1mm]
  % 空気を箱に分ける（横6・奥4・高さ3）
  \fill[bkBlue!6] \bkviaP(0,0,0) -- \bkviaP(6,0,0) -- \bkviaP(6,4,0) -- \bkviaP(6,4,3) -- \bkviaP(0,4,3) -- \bkviaP(0,0,3) -- cycle;
  \fill[bkOrange!30] \bkviaP(2,0,2) -- \bkviaP(3,0,2) -- \bkviaP(3,0,3) -- \bkviaP(2,0,3) -- cycle;
  \foreach \i in {0,...,6} {\draw[bkBlue!60!black, line width=0.4pt] \bkviaP(\i,0,0) -- \bkviaP(\i,0,3) -- \bkviaP(\i,4,3);}
  \foreach \k in {0,...,3} {\draw[bkBlue!60!black, line width=0.4pt] \bkviaP(0,0,\k) -- \bkviaP(6,0,\k) -- \bkviaP(6,4,\k);}
  \foreach \j in {0,...,4} {\draw[bkBlue!60!black, line width=0.4pt] \bkviaP(0,\j,3) -- \bkviaP(6,\j,3) -- \bkviaP(6,\j,0);}
  \draw[bkOrange!80!black, line width=1.2pt] \bkviaP(2,0,2) -- \bkviaP(3,0,2) -- \bkviaP(3,0,3) -- \bkviaP(2,0,3) -- cycle;
  % 風の矢印（上の面）
  \foreach \i/\j in {1/1, 3/2, 4.6/0.8, 2/3.1} {\draw[-{Stealth[length=3pt]}, bkRed, line width=1pt] \bkviaP(\i,\j,3) -- \bkviaP(\i+0.8,\j+0.5,3);}
  \bkviacloud{(30,42)}{1.1}
  \fill[bkGray!40] (0,0) -- (60,0) -- (60,-2) -- (0,-2) -- cycle;
  \node[anchor=north, font=\sffamily\footnotesize] at (30,-3) {地面};
  \draw[-{Stealth[length=3pt]}, bkInk, line width=0.7pt] (-5,0) -- (-5,24) node[anchor=south, font=\sffamily\footnotesize] {上空};
  % 説明
  \draw[bkOrange!80!black, line width=0.6pt] (30,20) -- (88,20);
  \node[anchor=west, draw=bkOrange!80!black, line width=0.8pt, fill=white, rounded corners=1.5mm, text width=38mm, align=left, font=\footnotesize, inner sep=2mm] at (88,20)
    {\textbf{1つの箱の空気の状態}\\気温・気圧（空気の重さでおす力）・風の向きと速さ・水蒸気の量};
\end{tikzpicture}
\bkfigcaption{空気を箱に分ける。1つの箱ごとに、気温や風などの「状態」を数で表す。}
\end{bkfigure}

まず、地球をおおう空気を、たて・横・高さに区切って、たくさんの箱に分けます。気象庁の計算では、地面を約\qty{13}{km}四方のます目に分け、高さの方向にも128の層に分けています。箱の数は、全部で数億こになります。

1つ1つの箱の「状態」は、気温・気圧・風の向きと速さ・水蒸気の量などの数です。そして、空気の動きには、物理の法則から作ったルールがあります。たとえば、次のようなルールです。
\begin{itemize}[leftmargin=1.2\zw, itemsep=1pt, topsep=2pt, label={\color{\bkTT}●}]
\item 空気は、気圧の高い所から低い所へおされて動く（風になる）。
\item 風は、温度や水蒸気を、となりの箱へ運ぶ。
\item あたたかい空気は軽くて上へ上がり、上がった空気は冷える。
\item 水蒸気をふくんだ空気が冷えると、雲や雨になる。
\end{itemize}
これらのルールを使うと、1つ1つの箱の「状態の変わる速さ」が、今の状態（その箱と、となりの箱の状態）から求められます。あとは、冷めるお湯と同じです。

\begin{taisetsu}
\centering 箱ごとに\enspace\textbf{次の状態}\enspace ＝\enspace\textbf{今の状態}\enspace ＋\enspace\textbf{変わる速さ}\enspace ×\enspace\textbf{数分}
\end{taisetsu}

1回の計算で、数分先へ進みます。これを何百回もくり返すと、明日の天気になります。だから、たくさんの計算を速くできるコンピュータ（スーパーコンピュータ）が必要なのです。

\Needspace{8\baselineskip}
\noindent\textbf{\bkhatenaref の答え}\par
天気予報は、明日の天気を1回の計算で求めているのではありません。いまの空気の状態を観測して、それを「はじめの状態」にします。そして、ルールを使って「数分先」を計算することを、何百回もくり返しているのです。ロケットや人工衛星の通り道も、同じ考え方で計算しています（④で見ます）。

「はじめの状態」は、地上の観測所（アメダス）、気象衛星「ひまわり」、気球につけた観測の機械（ラジオゾンデ）、船や飛行機などで、毎日はかっています。はじめの状態をできるだけ正しくはかることも、予報を当てるために大切です。


%% ---------------- ③ ボール・ロケット ----------------
\bksection{doguu}{ボールもロケットも、同じ考え方}

\bksubsection{5}{状態が2つの数のとき}

\label{bk:6-0:ball}投げたボールの動きも、同じ考え方で計算できます。ただし、ボールの「状態」は、1つの数では足りません。「どこにいるか（位置）」と「どれだけの速さで、どちらへ動いているか（速度）」の2つが必要です。

\begin{bkcenter}
\begin{bktable}{colspec={Q[l,m,wd=22mm] Q[l,m,wd=100mm]}, rowsep=2pt, row{1}={font=\sffamily\bfseries\small}}
求めるもの & ルール \\
次の位置 & 今の位置＋\textbf{今の速度}×$\Delta t$\quad （速度が、位置の変わる速さ） \\
次の速度 & 今の速度＋\textbf{今の加速度}×$\Delta t$\quad （加速度が、速度の変わる速さ） \\
\end{bktable}
\bktabcaption{ボールの状態（位置と速度）を、少しずつ進めるルール}
\end{bkcenter}

加速度は、ボールにはたらく力で決まります（運動の法則 $F=ma$、5-5）。空気の抵抗を考えないなら、ボールにはたらく力は重さ（重力）だけで、加速度は下向きにいつも\qty{9.8}{m/s^2}です。9.8は、地球の上で落とした物の速さが、1秒ごとに\qty{9.8}{m/s}ずつふえることを、実験ではかった値です。

\begin{bkfigure}
\begin{tikzpicture}[x=1mm, y=1mm]
  \fill[bkGreen!15] (-4,-2) rectangle (106,0);
  \draw[bkGreen!60!black, line width=0.8pt] (-4,0) -- (106,0);
  \bkkid{(-1,0)}{16}{bkBlue}{bkInk}
  % 0.1 秒ごとの位置（1 m を 10 mm、横は 6 m/s、たては 8 m/s で投げ上げ）
  \foreach \x/\y in {0/0, 6/7.51, 12/14.04, 18/19.59, 24/24.16, 30/27.75, 36/30.36, 42/31.99, 48/32.64, 54/32.31, 60/31.0, 66/28.71, 72/25.44, 78/21.19, 84/15.96, 90/9.75, 96/2.56} {
    \fill[bkOrange] ({\x+4},{\y+12}) circle (1.3);}
  \draw[-{Stealth[length=4pt]}, bkRed, line width=1pt] (4,12) -- (7.6,16.8);
  \node[anchor=west, font=\sffamily\footnotesize, text=bkRed!80!black, inner sep=1pt] at (9.5,13.6) {はじめの速度};
  \node[anchor=west, text width=28mm, align=left, font=\sffamily\footnotesize] at (102,30) {0.1秒ごとのボールの位置。点と点の間が広い所ほど、速く動いている。};
\end{tikzpicture}
\bkfigcaption{投げたボールの、0.1秒ごとの位置。はじめの位置と速度から、ルールで次々に計算できる。}
\end{bkfigure}

ゲームの中のボールやキャラクターも、この計算で動いています。画面は1秒に60回かきかえられるので、$\Delta t$ は $\dfrac{1}{60}$ 秒です（5-13のコラム）。人工衛星や、月や惑星へ向かう探査機の通り道（軌道）も、エンジンの力・地球や月が引く力などから加速度を求め、同じように少しずつ進めて計算します。

天気予報・冷めるお湯・ボール・人工衛星。どれも、「今の状態」と「次の瞬間を決めるルール」から、少しずつ未来へ進んでいます。

%% ---------------- ④ ルールとはじめの状態 ----------------
\bksection{doguu}{未来を決める2つのもの}

\bksubsection{6}{ルールと、はじめの状態}

\label{bk:6-0:futatsu}同じ部屋で、同じカップに、温度のちがうお湯を入れたら、どうなるでしょう。ルールは同じ「1分間に、部屋との差の5\,\%だけ変わる」です。ちがうのは、はじめの温度だけです。

\begin{bkfigure}
\begin{tikzpicture}
\begin{axis}[bkplot, width=104mm, height=56mm, xmin=0, xmax=66, ymin=0, ymax=92, xtick={0,20,40,60}, ytick={20,40,60,80},
  xlabel={時刻（分）}, ylabel={温度（℃）}]
  \addplot[bkthin, dashed, domain=0:62] {20};
  \addplot[bkcurve, domain=0:62, samples=60] {20+60*exp(-0.05*x)};
  \addplot[bkcurve2, domain=0:62, samples=60] {20+30*exp(-0.05*x)};
  \addplot[bkcurve3, domain=0:62, samples=60] {20-10*exp(-0.05*x)};
  \addplot[only marks, mark=*, mark size=1.8pt, bkInk] coordinates {(0,80) (0,50) (0,10)};
  \node[anchor=west, font=\sffamily\footnotesize, text=bkBlue] at (axis cs:5,84) {はじめ 80\,℃};
  \node[anchor=west, font=\sffamily\footnotesize, text=bkOrange!80!black] at (axis cs:5,55) {はじめ 50\,℃};
  \node[anchor=west, font=\sffamily\footnotesize, text=bkGreen!60!black] at (axis cs:4,6) {はじめ 10\,℃（冷たい水）};
  \node[anchor=west, font=\sffamily\footnotesize, text=bkGray, inner sep=1pt] at (axis cs:63,20) {部屋 20\,℃};
\end{axis}
\end{tikzpicture}
\bkfigcaption{同じルールでも、はじめの温度がちがうと、ちがう未来になる。どれも部屋の温度に近づく。}
\end{bkfigure}

80\,℃のお湯も、50\,℃のお湯も、冷めて20\,℃に近づきます。10\,℃の冷たい水は、反対にあたたまって20\,℃に近づきます。10\,℃のときは、部屋との差が $10-20=-10$ とマイナスになるので、「差の5\,\%だけ下がる」は「$10\times0.05=0.5$ だけ上がる」ことになるのです。

ルールだけでは、どのグラフになるか決まりません。はじめの状態も決めると、グラフはただ1本に決まります。5-13で、「法則とはじめの状態がそろって、はじめて先のようすが1つに決まる」と書いたのは、このことです。

\begin{taisetsu}
\centering\textbf{未来}\enspace は\enspace\textbf{ルール}\enspace と\enspace\textbf{はじめの状態}\enspace の2つで決まる
\end{taisetsu}

天気予報では、ルールは物理の法則で、はじめの状態は観測したデータです。では、はじめの状態を少しだけまちがえたら、どうなるでしょう。冷めるお湯なら、はじめが80\,℃でも80.1\,℃でも、ほとんど同じ未来になります。ところが空気の動きでは、はじめのごく小さなちがいが、日がたつにつれて大きく育つことがあります。そのため、天気をくわしく当てられるのは、数日から2週間ほど先までが限界だと考えられています。この話は、6-14のコラムで続きを読みます。

\begin{column}{天気を計算で予想する夢}
\noindent\begin{minipage}[t]{0.53\linewidth}
天気を計算で予想しようと考えた最初の一人は、イギリスのリチャードソンです。彼は、第一次世界大戦のころ、空気のルールを使って「6時間先の気圧」を手で計算しました。電卓もコンピュータもない時代で、計算には1か月以上かかりました。

ところが、出てきた答えは「6時間で気圧が約\qty{145}{hPa}も変わる」という、ありえない値でした。実際の気圧はほとんど変わっていなかったのです。のちの研究で、計算の方法がまちがっていたのではなく、はじめの状態のデータにあった気圧と風の小さなずれが、大きな見かけの変化を生んだことがわかっています。
\end{minipage}\hfill
\begin{minipage}[t]{0.43\linewidth}
リチャードソンは1922年に、このくわしい記録を本にしました。その中で、「6万4千人が大きな劇場に集まって、指揮者の合図で分担して計算すれば、天気の変化に計算が追いつく」という夢を書いています。

その夢は、コンピュータがかなえました。1950年、アメリカのコンピュータENIAC（エニアック）が、はじめて24時間先の予報を計算しました（計算にも約24時間かかりました）。日本の気象庁は、1959年にコンピュータでの予報を始めました。
\end{minipage}
\par\vspace{4pt}
{\centering
\begin{tikzpicture}[x=1mm, y=1mm]
  \draw[-{Stealth[length=4pt]}, bkNavy, line width=1.1pt] (0,0) -- (118,0);
  \foreach \x/\y/\t in {6/1922/{リチャードソン\\の本（手計算）}, 38/1950/{ENIAC で\\24時間先の予報}, 70/1959/{気象庁が\\コンピュータの予報}, 102/いま/{スーパーコンピュータ\\約13 km のます目}} {
    \fill[bkNavy] (\x,0) circle (1.2);
    \node[anchor=south, font=\sffamily\bfseries\footnotesize, text=bkNavy] at (\x,2) {\y};
    \node[anchor=north, font=\sffamily\footnotesize, align=center, inner sep=0pt] at (\x,-2.5) {\t};}
\end{tikzpicture}\par}
\end{column}

%% ---------------- ⑤ この巻の地図 ----------------
\bksection{goal}{この巻でできるようになること}

\label{bk:6-0:chizu}この巻では、「次の瞬間を決めるルール」を式で書いて、未来を計算します。変わる速さは微分（第5巻）で書けるので、このルールの式は、微分の入った方程式になります。これを\textbf{微分方程式}（びぶんほうていしき）といいます（6-1でくわしく）。\bkidx{びぶんほうていしき}{微分方程式}

第6巻は、3つの部に分かれています。

\begin{bkfigure}
\begin{tikzpicture}[x=1mm, y=1mm]
  \bkviapart{0}{第1部\enspace 微分方程式の入口}{6-1〜6-7}%
    {\draw[bkInk, line width=0.5pt] (-14,-8) -- (14,-8) (-14,-8) -- (-14,9);
     \draw[bkBlue, line width=1.1pt, domain=0:26, samples=30] plot ({-14+\x}, {-8+16*exp(-\x/9)});
     \foreach \k in {0,...,5} {\fill[bkOrange] ({-14+4*\k},{-8+16*(1-4/9)^\k}) circle (0.9);}}
    {\bkviaitem{ルールを微分方程式で書いて、読む。}\bkviaitem{表計算で少しずつ進めて解く（オイラー法）。}\bkviaitem{積分だけで解ける形と、答えが指数関数になる形。}\bkviaitem{雨つぶ・ばね・ゆれながら止まるばね。}}
  \bkviapart{-42}{第2部\enspace 複素数と双曲線関数}{6-8〜6-11}%
    {\draw[bkInk, line width=0.5pt] (-13,0) -- (13,0) (0,-10) -- (0,10);
     \draw[bkInk!70, line width=0.6pt] (0,0) circle (9);
     \draw[-{Stealth[length=3pt]}, bkPurple, line width=1.2pt] (0,0) -- (40:9);
     \draw[-{Stealth[length=3pt]}, bkOrange, line width=1.2pt] (0,0) -- (130:9);}
    {\bkviaitem{2乗すると $-1$ になる新しい数 $j$（虚数）。}\bkviaitem{かけ算を、平面の上の「回転」として見る。}\bkviaitem{オイラーの公式で、指数関数と三角関数をつなぐ。}\bkviaitem{電線のたるみの形（双曲線関数）。}}
  \bkviapart{-84}{第3部\enspace 式で解く}{6-12〜6-15}%
    {\draw[bkInk, line width=0.5pt] (-14,0) -- (16,0) (-14,-9.5) -- (-14,9.5);
     \draw[bkBlue, line width=1.1pt, domain=0:30, samples=80] plot ({-14+\x}, {8.5*exp(-\x/10)*cos(\x*40)});}
    {\bkviaitem{ゆれながら止まる動きを、2次方程式で解く。}\bkviaitem{ゆすられて大きくゆれる「共振」。}\bkviaitem{いくつもつながった物の動き。}\bkviaitem{ばね・熱・電気の式を、くらべてまとめる。}}
\end{tikzpicture}
\bkfigcaption{第6巻の3つの部と、できるようになること}
\end{bkfigure}

\bksubsection{7}{2周読む}

\label{bk:6-0:nishu}この巻は、同じ場面を2回通ります。1周めの第1部では、まず表計算で数とグラフを出して、答えのようすを先に見ます（1階の式は、第1部の6-3〜6-5で式でも解きます）。2周めの第3部では、ゆれる動きなどを、第2部で手に入れた新しい数（複素数）を使って式で求めます。そして、式の答えのグラフを、1周めの表計算のグラフに重ねて、「ぴったり重なった」ことを確かめます。

\begin{bkfigure}
\begin{tikzpicture}[x=1mm, y=1mm,
  pb/.style={bk node, minimum width=36mm, minimum height=16mm, font=\sffamily\small}]
  \node[pb] (p1) at (0,0) {\textbf{1周め}（第1部）\\{\footnotesize 表計算で、数とグラフ}};
  \node[pb] (p2) at (46,0) {\textbf{道具}（第2部）\\{\footnotesize 複素数・オイラーの公式}};
  \node[pb] (p3) at (92,0) {\textbf{2周め}（第3部）\\{\footnotesize 式で解いて、重ねる}};
  \draw[bk link] (p1) -- (p2);
  \draw[bk link] (p2) -- (p3);
  \draw[bk link, bkOrange!80!black] (p3.south) -- ++(0,-8) -| node[bk link label, pos=0.25, below, text=bkOrange!80!black] {式の答えを、1周めのグラフに重ねて確かめる} (p1.south);
\end{tikzpicture}
\bkfigcaption{第6巻は2周読む。先に数で見て、あとで式で確かめる。}
\end{bkfigure}

第3部は、大学で学ぶ内容です。1回読んでわからなくても、だいじょうぶです。第1部の表計算だけでも、未来はちゃんと計算できます。式の答えは、「表計算で見た答えが、なぜそうなるのか」を教えてくれる、2周めのおまけだと思って読み進めましょう。

%% ---------------- ⑥ 3つの分野 ----------------
\bksubsection{8}{ばね・熱・電気回路が、くり返し出てくる}

\label{bk:6-0:bunya}この巻では、3つの分野の場面が、何回も出てきます。機械の「ばね・質量（おもり）・ダンパー」、「熱」、そして「電気回路」です。

\begin{bkfigure}
\begin{tikzpicture}[x=1mm, y=1mm]
  \useasboundingbox (-4,-9) rectangle (132,31);
  % --- ばね・質量・ダンパー ---
  \begin{scope}[shift={(0,0)}]
    \fill[bkGray!25] (-3,-1.6) rectangle (40,0);
    \draw[bkGray!80!black, line width=0.8pt] (-3,0) -- (40,0);
    \bkviawall{(0,0)}{27}
    \bkviaspring{(0,20)}{22}{4}{2}{bkInk}
    % ダンパー（油の入った筒と、ピストン）
    \draw[bkInk, line width=0.9pt] (0,10) -- (7,10);
    \fill[bkBlue!20] (7,7) rectangle (16,13);
    \draw[bkInk, line width=0.9pt] (16,7) -- (7,7) -- (7,13) -- (16,13);
    \draw[bkInk, line width=1.2pt] (12,8) -- (12,12);
    \draw[bkInk, line width=0.9pt] (12,10) -- (22,10);
    % 質量
    \filldraw[fill=bkOrange!40, draw=bkOrange!70!black, line width=0.9pt] (22,2.4) rectangle (36,25);
    \fill[bkInk] (25.5,1.2) circle (1.2) (32.5,1.2) circle (1.2);
    \node[font=\sffamily\footnotesize] at (29,14) {$m$};
    \node[anchor=south, font=\sffamily\footnotesize] at (11,22.5) {ばね $k$};
    \node[anchor=north, font=\sffamily\footnotesize] at (11.5,5.6) {ダンパー $c$};
    \node[anchor=north, font=\sffamily\bfseries\footnotesize, text=bkDeep] at (18,-4) {ばね・質量・ダンパー};
  \end{scope}
  % --- 熱 ---
  \begin{scope}[shift={(50,0)}]
    \filldraw[fill=bkOrange!12, draw=bkInk!70, line width=0.9pt] (2,0) -- (0.5,19) -- (21.5,19) -- (20,0) -- cycle;
    \fill[bkOrange!35] (1.9,1.2) -- (0.8,15) -- (21.2,15) -- (20.1,1.2) -- cycle;
    \draw[bkInk!70, line width=0.9pt] (21.1,14) arc (90:-90:3.8 and 4.6);
    \foreach \s in {5,11} {\draw[bkGray, line width=0.7pt] (\s,21) .. controls (\s-1.5,23.5) and (\s+1.5,25) .. (\s,27.5);}
    % 温度計
    \draw[bkInk, line width=0.7pt, rounded corners=0.8mm] (15.2,8) rectangle (17.8,29);
    \fill[bkRed] (16.5,7.8) circle (1.8);
    \fill[bkRed] (16.1,8) rectangle (16.9,22);
    \node[anchor=north, font=\sffamily\bfseries\footnotesize, text=bkDeep] at (11,-4) {熱（冷めるお湯）};
  \end{scope}
  % --- 電気回路 ---
  \begin{scope}[shift={(90,0)}]
    % 電池（左）・スイッチ（上）・抵抗（右）・コンデンサー（下）
    \draw[bkInk, line width=0.9pt] (0,2) -- (0,9) (0,15) -- (0,24) -- (8,24);
    \draw[bkInk, line width=0.9pt] (8,24) -- (15,27.5);
    \fill[bkInk] (8,24) circle (0.8) (16,24) circle (0.8);
    \draw[bkInk, line width=0.9pt] (16,24) -- (36,24) -- (36,19);
    \draw[bkInk, line width=0.9pt, fill=white] (34,9) rectangle (38,19);
    \draw[bkInk, line width=0.9pt] (36,9) -- (36,2) -- (21,2);
    \draw[bkInk, line width=1.4pt] (21,-1.5) -- (21,5.5) (17.5,-1.5) -- (17.5,5.5);
    \draw[bkInk, line width=0.9pt] (17.5,2) -- (0,2);
    \draw[bkInk, line width=1.4pt] (-3.5,15) -- (3.5,15);
    \draw[bkInk, line width=2.6pt] (-2,9) -- (2,9);
    \node[anchor=west, font=\sffamily\footnotesize] at (3.5,12) {電池};
    \node[anchor=south, font=\sffamily\footnotesize] at (12,28) {スイッチ};
    \node[anchor=east, font=\sffamily\footnotesize] at (32.5,14) {抵抗 $R$};
    \node[anchor=south, font=\sffamily\footnotesize, fill=white, inner sep=0.8pt] at (19.25,7) {$C$};
    \node[anchor=north, font=\sffamily\bfseries\footnotesize, text=bkDeep] at (18,-4) {電気回路};
  \end{scope}
\end{tikzpicture}
\bkfigcaption{この巻でくり返し出てくる3つの場面。ばね・ダンパーの $k$・$c$・$m$、回路の $R$・$C$ は、それぞれの章で説明する。}
\end{bkfigure}

\begin{bkcenter}
\begin{bktable}{colspec={Q[c,m,wd=10mm] Q[l,m,wd=40mm] Q[l,m,wd=32mm] Q[l,m,wd=40mm]}, rowsep=0pt, colsep=4pt, cells={font=\small}, row{1}={font=\sffamily\bfseries\small}}
章 & ばね・質量・ダンパー & 熱 & 電気回路 \\
6-1 & ばねのルールを読む & 冷めるお湯 & コンデンサーの放電 \\
6-2 & ばねを表計算で動かす & 80\,℃のお湯 & 充電を表計算で \\
6-3 & 落ちるボール & ヒーターで温める & 一定の電流で充電 \\
6-4 & --- & 冷める・温まる & 充電・放電、コイル \\
6-5 & 雨つぶ・エレベーター & --- & コイルの電流 \\
6-6 & ばねと振り子 & --- & 電気の振動 \\
6-7 & ゆれながら止まる & --- & ゆれながら止まる電気 \\
6-12 & ダンパーつきのばね & --- & ちょうどよい抵抗 \\
6-13 & ゆすられるばね（共振） & --- & ラジオの同調 \\
6-14 & 2つのおもり & 2つの物の温度 & 2段の回路 \\
6-15 & \SetCell[c=3]{l} 3つの分野が、同じ形の式で書けることをまとめる & & \\
\end{bktable}
\bktabcaption{3つの分野の場面が出てくる章}
\end{bkcenter}

電気回路では、とくに「\textbf{スイッチを入れた瞬間}」からのようすを調べます。スイッチを入れた瞬間、電気の量はすぐには決まった値にならず、少しの時間をかけて変わっていきます。この「変わっていくとちゅう」のようすを、\textbf{過渡応答}（かとおうとう）といいます。\bkidx{かとおうとう}{過渡応答}電気製品の電源を入れたときにも、回路の中ではいつも過渡応答が起きています。


ちがう分野の場面なのに、同じ形の式で書けることが何回も出てきます。1つの場面で答えの形がわかれば、ほかの分野にもそのまま使えます。それが、数学で考えることの大きな強みです。

%% ---------------- ⑦ 乗りかえ駅 ----------------
\bksection{bkviaroute}{この巻が通る路線と乗りかえ駅}

\label{bk:6-0:rosen}この本の章は、7本の路線の駅です。第6巻を通る5本の路線のうち、「変化」「たし合わせ」「倍々」「回転と波」の4本は、この巻の駅（6-1・6-3・6-4・6-10）で乗りかえて、物理と工学の終点へ向かいます。

\begin{bkfigure}
\begin{tikzpicture}[x=1mm, y=1mm]
  \bkviaroute{56}{1}{5-1}{{\bfseries 6-1}・6-2・6-5・6-15}{8-4}
  \bkviaroute{42}{2}{5-10}{{\bfseries 6-3}}{8-11}
  \bkviaroute{28}{3}{5-14}{{\bfseries 6-4}・6-11}{9-5}
  \bkviaroute{14}{4}{5-15}{6-6〜6-9・{\bfseries 6-10}・6-12・6-13}{9-1}
  \bkviaroute{0}{5}{3-6}{{\bfseries 6-14}}{7-1}
  \node[font=\sffamily\bfseries\small, text=bkGray] at (24,66) {前の巻};
  \node[font=\sffamily\bfseries\small, text=bkDeep] at (74,66) {この巻（第6巻）};
  \node[font=\sffamily\bfseries\small, text=bkGray] at (124,66) {この先};
\end{tikzpicture}
\bkfigcaption{第6巻を通る5本の路線。四角の中がこの巻の駅（太字は大きな駅）。}
\end{bkfigure}

\begin{bkcenter}
\begin{bktable}{colspec={Q[c,m,wd=7mm] Q[l,m,wd=20mm] Q[l,m,wd=92mm]}, rowsep=0.8pt, cells={font=\small}, row{1}={font=\sffamily\bfseries\small}}
 & 路線 & 大きな駅のつながり \\
\bkrouteicon{1} & 変化 & ひき算（1-5）→ 傾き（3-5）→ 微分（5-1）→ \textbf{微分方程式}（6-1） \\
\bkrouteicon{2} & たし合わせ & かけ算の面積（2-3）→ 積分（5-10）→ \textbf{積分するだけで解ける形}（6-3） \\
\bkrouteicon{3} & 倍々 & 指数関数（3-11）→ $e$（5-14）→ \textbf{自分に比例してへる}（6-4） \\
\bkrouteicon{4} & 回転と波 & 三角関数（4-8）→ $\sin$・$\cos$ の微分（5-15）→ \textbf{オイラーの公式}（6-10） \\
\bkrouteicon{5} & 表 & 表計算（2-18）→ 連立方程式（3-6）→ \textbf{連立微分方程式}（6-14） \\
\end{bktable}
\bktabcaption{5本の路線の、この巻の大きな駅}
\end{bkcenter}

変化の路線（\bkrouteicon[0.9em]{1}）は、この巻の中心です。「ひき算で変化を見る」から始まった道が、6-1で「変化のルールから未来を求める」微分方程式に着きます。回転と波の路線（\bkrouteicon[0.9em]{4}）は、この巻でいちばん多くの駅に止まります。ばねのゆれも、電気のゆれも、回る矢印の影として見ることができるからです（6-10）。

%% ---------------- ⑧ 新しい記号 ----------------
\bksection{card}{この巻で出会う新しい記号}

\label{bk:6-0:kigo}第6巻では、次のような新しい記号に出会います。いまは意味がわからなくてもだいじょうぶです。どの記号も、使う前に場面と図でていねいに説明します（記号カードは巻末にも並べてあります）。

\begin{bkcenter}
\begin{bktable}{colspec={Q[c,m,wd=40mm] Q[l,m,wd=28mm] Q[l,m,wd=44mm] Q[c,m,wd=9mm]}, rowsep=0.3pt, colsep=3pt, cells={font=\small}, row{1}={font=\sffamily\bfseries\small}}
記号 & 読み方 & ひとことで言うと & 章 \\
$\dfrac{\dd y}{\dd t}=f(y)$ & ディーワイ・ディーティー イコール エフ・オブ・ワイ & 変わる速さを、今の状態で決めるルール（微分方程式） & 6-1 \\
$y_{n+1}=y_n+f(y_n)\Delta t$ & （オイラー法の1行） & 今から、少し先へ進む計算 & 6-2 \\
$\tau$ & タウ & 時定数。変化の速さの目安 & 6-4 \\
$v_{\mathrm f}$、$v_{\mathrm t}$ & ブイ・エフ、ブイ・ティー & 終端速度（それ以上速くならない速さ） & 6-5 \\
$\omega_{\mathrm n}$ & オメガ・エヌ & 固有角振動数（その物が自分でゆれる速さ） & 6-6 \\
$\zeta$ & ゼータ & 減衰比（ブレーキのきき方） & 6-7 \\
$j$ & ジェイ & 2乗すると $-1$ になる数（虚数単位） & 6-8 \\
$\lvert z\rvert$、$\arg z$ & ゼットの絶対値、偏角 & 複素数の大きさと向き & 6-9 \\
$e^{j\theta}$ & イーのジェイシータ乗 & 回る矢印（オイラーの公式） & 6-10 \\
$\cosh$、$\sinh$、$\tanh$ & ハイパボリック・コサインなど & 双曲線関数 & 6-11 \\
$ms^2+cs+k=0$ & （特性方程式） & ゆれ方を決める2次方程式 & 6-12 \\
\end{bktable}
\bktabcaption{この巻で出会う記号（くわしくはそれぞれの章で）}
\end{bkcenter}

電気回路の部品の量は、記号のあとに単位を〔　〕でくくって書きます。たとえば、抵抗 $R$〔$\Omega$〕、コンデンサーの静電容量 $C$〔F〕、コイルのインダクタンス $L$〔H〕、電圧 $V$〔V〕、電流 $i$〔A〕です。記号の $C$（静電容量）と、電気の量の単位の C（クーロン）のように、同じ字が別の意味で出てくることがあるからです。

%% ---------------- この巻の読み方 ----------------
\bksection{hon}{この巻の読み方}

\label{bk:6-0:yomikata}
\begin{bkfigure}
\begin{tikzpicture}[x=1mm, y=1mm]
  \bkviatalk{0}{sora}{ソラ}{第6巻からは、大学で学ぶ内容が多いんだって。微分方程式って、名前がもうむずかしそう……。}
  \bkviatalk{-24}{delta}{デルタ}{だいじょうぶ。この巻のルールは、冷めるお湯の「1分で差の5\,\%下がる」と同じ考え方だよ。先に表計算で数とグラフを見て、あとから式で確かめるよ。1回で分からなくても、第1部の表計算ができれば、未来は計算できるよ。}
\end{tikzpicture}
\end{bkfigure}

\begin{bkfigure}
\begin{tikzpicture}[x=1mm, y=1mm]
  \bkviacard{(0,0)}{shiki}{bkNavy}{①\enspace 式の変形は1行ずつ}
    {式の変形は、1行に1つのことだけをします。何をしたかを右に書きます。\par\vspace{2pt}
     \bkviaitem{「両辺を2でわる」「$e^0=1$ を使う」のように書く。}\bkviaitem{わからない行は、1つ前の行とくらべてみよう。}}
  \bkviacard{(68,0)}{te}{bkGreen}{②\enspace 表計算とプログラム}
    {表計算（2-18）とプログラム（2-19）を、たくさん使います。\par\vspace{2pt}
     \bkviaitem{表計算の画面と、Python のプログラムの枠をのせてあります。}\bkviaitem{自分のコンピュータで、数を変えてためしてみよう。}}
  \bkviacard{(0,-57)}{check}{bkOrange}{③\enspace 単位で確かめる}
    {答えが出たら、単位を確かめます（2-10）。\par\vspace{2pt}
     \bkviaitem{「速さ」なら ℃/分、V/s のように「÷時間」の単位になるはず。}\bkviaitem{換算は、\qty{60}{s/min} のような単位つきの1を使う。}}
  \bkviacard{(68,-57)}{goal}{bkDeep}{④\enspace ★5の章は3段}
    {大学の内容の章は、3つの段に分けてあります。\par\vspace{2pt}
     \bkviaitem{「目でわかる」「手でわかる」「式でわかる」。}\bkviaitem{「ここまでで十分」の旗まで読めば、その章のゴールです。}}
\end{tikzpicture}
\end{bkfigure}
